% Part: counterfactuals
% Chapter: minimal-change-semantics
% Section: transitivity

\documentclass[../../../include/open-logic-section]{subfiles}

\begin{document}

\olfileid{cnt}{min}{cpo}

\olsection{प्रतिपरिवर्तन}

वास्तविक आणि कठोर
अभिव्यंजने त्यांच्या
प्रतिपरिवर्तनांशी सममूल्य
असतात. प्रतिवास्तविक
विधानांच्या बाबतीत
तसे नाही. क्रात्सर
यांनी दिलेले उदाहरण
पाहा:
\begin{quote}
  गटे यांचा मृत्यू १८३२ मध्ये
  झाला नसता, तरी आता
  ते मृतच असते.

  गटे आता मृत नसते,
  तर त्यांचा मृत्यू
  १८३२ मध्ये झाला असता.
\end{quote}
पहिले वाक्य सत्य आहे:
माणसे शेकडो वर्षे
जगत नाहीत. दुसरे
स्पष्टपणे असत्य आहे:
गटे आता मृत नसते,
तर ते अजून जिवंत
असते; त्यामुळे
त्यांचा मृत्यू
१८३२ मध्ये
झालेला नसता.

\begin{ex}\ollabel{ex:contraposition-counterex}
  गोलक चिन्हार्थमीमांसेत
  प्रतिपरिवर्तन अवैध ठरते;
  म्हणजे
  $p \cif q \Entails/
  \lnot q \cif \lnot p$.
  $p$ चा अर्थ ‘‘गटे
  यांचा मृत्यू १८३२
  मध्ये झाला नाही’’
  आणि~$q$ चा अर्थ
  ‘‘गटे आता मृत
  आहेत’’ असा घ्या.
  हे पुढील प्रतिमानात
  दाखवता येते:
  $\mModel{M_1} = \tuple{W, O, V}$,
  येथे $W =
  \{w, w_1, w_2\}$,
  $O_w = \{\{w\}, \{w, w_1\},
  \{w, w_1, w_2\}\}$,
  $V(p) = \{w_1, w_2\}$
  आणि $V(q) = \{w, w_1\}$.
  $O_{w_1}=\{\{w_1\}\}$ आणि $O_{w_2}=\{\{w_2\}\}$ घ्या;
  $O$ आता संपूर्ण $W$ वर परिभाषित आहे.
  म्हणून~$w$ हे वास्तविक
  जग आहे, ज्यात गटे
  यांचा मृत्यू १८३२
  मध्ये झाला आणि
  ते अजूनही मृत आहेत;
  $w_1$ हे (जवळचे)
  जग आहे, ज्यात
  त्यांचा मृत्यू,
  समजा, १८३३ मध्ये
  झाला आणि ते
  अजूनही मृत आहेत;
  तर~$w_2$ हे
  (दूरचे) जग आहे,
  ज्यात गटे
  अजून जिवंत आहेत.
\begin{figure}
\begin{center}
\begin{tikzpicture}[scale=.6,layerwidth=1.5]\tiny
  \spheresystem{3}
  \begin{scope}
    \clip \propositionplot{0}{.8}{50}{4.5} -- (4.5,-4.5) -- (-4.5,-4.5) -- (-4.5,4.5) -- (4.5,4.5);
    \spherelayer{3}
  \end{scope}
  \propositionintersect{0}{2}{110}{5.5}
  \proposition{0}{.8}{50}{4.5}
  \path (0,0) node[world] {$w$};
  \spherepos{0}{2}{node[world] {$w_1$}}
  \spherepos{-50}{3}{node[world] {$w_2$}}
  \spherepos{28}{3}{node {$q$}}
  \spherepos{42}{3}{node {$\lnot q$}}
  \spherepos{-55}{4}{node {$p$}}
  \spherepos{-67}{4}{node {$\lnot p$}}
\end{tikzpicture}
\caption{प्रतिपरिवर्तनाचे प्रतिउदाहरण}
\ollabel{fig:contraposition}
\end{center}
\end{figure}
$p$ समाविष्ट करणारा
$S = \{w, w_1\}$
हा गोलक आहे आणि
त्यातील सर्व जगांत
$p \lif q$ सत्य आहे;
म्हणून
$\mSat{M_1}{p \cif q}[w]$.
पण $\lnot q$
समाविष्ट करणाऱ्या
$\{w, w_1, w_2\}$
गोलकात~$w_2$
हे जग आहे; तेथे
$q$ असत्य आणि
$p$ सत्य आहे.
म्हणून
$\mSat/{M_1}{\lnot q
\lif \lnot p}[w_2]$. हा $O_w$ मधील एकमेव $\lnot q$-जग
धारण करणारा गोलक असल्याने
$\mSat/{M_1}{\lnot q\cif\lnot p}[w]$.
मूळ प्रतिवास्तविक विधान सत्य असून त्याचे प्रतिपरिवर्तन असत्य आहे.
\end{ex}

\end{document}
